\emph{By Steve Leigh}

There are a bewildering variety of revolutionary left groups both in the U.S. and internationally. This article is an attempt to explain the relation of Red Spark to the rest of the revolutionary Left. It will look historically at the splits in the revolutionary movement and explain where Red Spark stands on each of those divisive issues. Also, even though a political group was at one time on one side or the other of a split, they may move to the other side in a different historical period. The same issues often re-arise in different forms later. For example, some groups that started out as revolutionary later moved toward reformism while still using revolutionary language. Some that came from the Trotskyist tradition later moved toward Stalinism.

\paragraph{Scientific vs. Utopian Socialism}

The first major split was between the Marxists, or scientific socialists, and the Utopians in the early 1800's. For Marx, socialism was the self-emancipation of the working class. He saw history as a pattern of class struggle based on the roles of different social classes in the economy---in the ``mode of production''. Socialism for Marx was a definite historical stage that would succeed capitalism after the working class organized to take political and economic control of society. History was heading toward socialism, but socialism was not inevitable. The working class would have to take power in an organized, conscious revolution. Marx and Engels saw socialism and communism as synonymous.

Opposed to Marx were the Utopians who did not base socialism on class struggle or even on class. They instead believed in educating all of society---or even just the ruling class---to the superiority of their particular Utopia. A common feature of these Utopias was a top-down approach. It would be instituted by convincing the rich to fund it or the legislators to pass laws. This was ``socialism from above''(as noted by Hal Draper). It denied the essential element in Marx's socialism ``the self-activity of the working class''. This division over democracy vs. elitism and revolution vs. reform was fundamental. It reappeared in the split between social democracy and revolutionary socialism within what was in theory a Marxist world party, the Second international, after 1914. An example of Utopian socialists are those who try to ``go back to the land'' and set up communes separate from society. They hope that the world will change when influenced by their example, again a purely educational approach. They are idealists who see ideas as central. Marxists on the other hand are materialists who see ideas as flowing from material conditions. Changing material conditions is the important thing. Convincing people of ideas is part of the process of changing material conditions.

Another early split was between Marx and the Anarchists. Anarchists thought that political authority was the root of exploitation and oppression. They were thus opposed to having even a revolutionary workers' state after a workers' insurrection. Marx called for a ``dictatorship of the proletariat'', or a workers' democracy. This state would be necessary to crush the resistance of the capitalists and to make the concerted economic changes needed to bring about the transition to socialism. Anarchists thought that this new workers' state would lead to new oppression. For Marx, the dictatorship of the proletariat was a transitional phase until all classes were abolished. Once full communism arrived (no money, no classes, no division between menial and manual labor) there would be no need for any government. For Marx, as long as the revolutionary government was a real workers' democracy, there would be little danger of it degenerating into a new tyranny. Without such a government, the transition to communism would be sabotaged by the capitalists. Another major point on which Marx and the Anarchists differed was on the need for a political party. For Marx an independent political party of the working class was needed to counter the political organization of the capitalists. Anarchists saw the party as another form of authority and hence oppressive. Firebrand is on Marx's side of these early disputes.

\paragraph{Revolutionaries vs. Reformists}

The next major split in the socialist tradition was between the reformists (the social democrats) and the revolutionaries. In 1917, the Russian working class seized power, destroyed the capitalist state, and created a worker's state, as Marx had predicted. They nationalized industry under workers' control and suppressed the capitalists. At the same time the mainstream Socialist parties of the Second International rejected this road to socialism. They took Marx's emphasis on political action to mean that electing a majority to Congress or the Parliament was enough to abolish capitalism. They felt that the capitalist state machine (courts, army, police, bureaucracy) could be taken over and used by the workers to create socialism. They soon came to feel that capitalism could be reformed a little at a time into socialism by passing laws and hence elections were more important than class struggle. Because they wanted to take over the capitalist state, not smash it, they had to prove their reliability and trustworthiness to the capitalists. This meant supporting their own governments in World War I and actually opposing revolutionary challenges to capitalism. In order to reform capitalism, they had to preserve it. As Marx predicted, this attempt to reform capitalism slowly into socialism instead of overthrowing it has never worked. After over 100 years of Social Democratic governments in Europe, capitalism is as strong there as ever. Reformism is a revival of the top-down approach that Marx rejected in his split with the Utopians. It meant relying on legislators rather than workers to change society. It failed to realize that relations of production are the fundamental social relations. These relations can only be altered by the people involved in those relations, the workers themselves. Reformism was also based on fatalism. It saw socialism as the inevitable result of capitalism. This fatalism fit well with its gradualism. Fatalism denied the need for concerted deliberate action, much less revolution.

The social democratic/reformist tradition is represented by all the established Socialist and Labor parties in Europe. In the U.S. most of the Democratic Socialists of America (D.S.A.) accepts this overall approach. Some in D.S.A. attempt to reform the Democratic Party, which is a capitalist party, into a Labor Party, which would then attempt to reform capitalism into socialism. This is called Realignment. Others would like to use ``socialist'' candidates and office holders in the Democratic Party to create the basis to build a new party. This is called the ``dirty break'' strategy. Even this, slightly more left-wing strategy is still reformist. The goal of the new party would be to use the capitalist state to pass legislation which would usher in socialism, along with mass action from below. Those who seek rapid transition through the capitalist state look to Karl Kautsky, one of the main leaders of the German Social Democratic Party in the early 1900s. They call their strategy Democratic Socialism which they try to differentiate from gradualist reformist Social Democracy.

Even the social-democratic parties in Europe (Socialist Party of France, Labor Party of Britain, Social Democrats of Germany etc.) have dropped the formal goal of socialism from their programs. All that is left is the continual reform of capitalism and in reality not much of that. Social Democrats today have often embraced neoliberalism (the idea that only the market can be a true regulator of the economy).

Though no other socialist group in the U.S. is as clearly reformist as D.S.A., many ``revolutionary'' groups have moved back toward reformism at least to some extent. The Communist Parties, while still claiming a Leninist basis, have become almost indistinguishable from the reformists. In the U.S., they support the Democratic Party. Communist Parties have supported or joined countless Socialist Party governments and even formed Popular Front coalitions with capitalist parties from the 30's on.

Even some Trotskyist groups have moved in this direction. The Militant Tendency in Britain entered the Labor Party in the 1940's to meet workers and ended up justifying its entry by rejecting the need for an independent revolutionary party. It even modified the Marxist and Leninist idea of the need to smash the state. Instead, it called for electing a Labor government with a ``socialist program''. It also softened its approach to imperialism. An offshoot of the Militant Tendency, the Committee for a Workers' International, still holds many of these politics. Socialist Alternative in the U.S. and its worldwide comrades, International Socialist Alternative comes from this tradition. It identifies socialism with a list of demands that reform the worst aspects of capitalism, instead of calling for a workers' state to replace the capitalist state. It accepts as a basic programmatic statement an article called ``Marxism and the State'' which says that socialism can be achieved in two ways,``either by workers' revolution or by the capitalist state nationalizing the economy''. On international issues, it refused to support the Iraqi resistance against the U.S., and calls for ``socialist apartheid'' in Palestine. It calls for two separate ``socialist'' states, one for Jews and one for Palestinians. It does not call for the smashing of the Zionist state of Israel. In the Gaza Crisis of 2023--4 it condemned Hamas' military attack on Israel. Softness on the capitalist state often leads to softness on imperialism.

\paragraph{The Bolshevik Tradition}

The other side of this fundamental division with the Reformist/Social-Democratic/Socialism-from-Above tendency was the continuation of Marxism represented by the revolutionary Bolshevik Party in Russia led by Lenin. It opposed World War I and after its own successful revolution in 1917 it formed the Communist International in 1919 to encourage revolution throughout the world. Lenin knew that only by having a revolutionary vanguard party to lead a workers' movement would it be able to resist the highly organized reformist and capitalist parties. For Lenin the vanguard was not an elite group outside the working class, but was simply the most militant class-conscious wing of the workers' movement. Its purpose was to convince the working class as a whole to take power, abolish the capitalist state and capitalism itself. Firebrand of course identifies with this Leninist tradition.

\paragraph{Ultra-Leftism}

Within the revolutionary wing another split occurred around 1920. Lenin and the Bolshevik Party rejected reformism but not the fight for reforms. They felt that in non-revolutionary periods, the revolutionary party should ally with reformist workers to win partial demands. In this alliance, which was called a United Front, revolutionaries would maintain their independent organization and political perspective. They would try to convince the other workers to become revolutionaries by propaganda about the long-term goal and by proposing more effective methods of struggle for the common goals. This meant ``going where the workers are''---i.e. becoming members of unions even if these unions were led by reformists or even reactionaries. The problem for revolutionaries was to do two things: preserve their politics and relate those politics to non-revolutionary workers. To go too far in either direction was to quit being an effective revolutionary. One extreme is to hide Marxist politics for fear of alienating moderates. The other extreme is to be afraid of diluting revolutionary politics and refusing to work with people who disagree even when there are other points of agreement to work together on.

In the Communist International some groups veered in the latter direction. They refused to work in unions led by non-revolutionaries. Instead, they tried to form pure ``red unions'' made up only of communists and their sympathizers. Also, on principle they refused to run candidates in bourgeois elections for educational purposes as the Bolsheviks had done. Of course, for Lenin, socialism could not come by electing socialists to office but elections could be a useful way to spread ideas and organization. He would never of course support capitalist candidates. Pure examples of what Lenin called ``infantile ultra-leftism'' are few in number today.

One group that explicitly follows the ultra-left line is the International Communist Current. Some groups that claim to be in the Leninist tradition actually are closer to ultra-leftism. One example of this is the Spartacist League which expresses this tendency in its sectarian orientation to savagely denouncing other groups in its press to the near exclusion of general political analysis. Another group with a similar approach is the Socialist Equality Party which produces the World Socialist Website. Another example of this is the Maoist Revolutionary Communist Party (RCP). It shows the elitism that often accompanies ultra-leftism by lack of support of most economic struggles by workers especially if organized by the established unions. They see most of the American working class as ``bought off'' and incapable of revolution. They favor a minority revolution of the most committed as shown by their support of Sendero Luminoso in Peru and Mao in China. This is another form of socialism from above: when the majority of the working class is too conservative to make a revolution, the revolutionaries simply impose it on them. This is a far cry from the ``self-emancipation of the working class'' or Marx or Lenin's approach of ``patiently explaining'' revolutionary ideas until a majority were convinced. For Marxists, a revolutionary situation cannot be created by the will of a party. Though a party IS absolutely essential in a revolutionary situation to ensure victory to the workers, only objective material conditions can create the situation. Ultra-leftists thus err in the direction of voluntarism instead of the fatalism of the reformists and idealism instead of the mechanical determinism of social-democrats.

\paragraph{Stalinism vs. Leninism}

Within the Leninist tradition, a new split soon developed. After the seizure of power by the workers in 1917, the Russian Revolution was attacked by 14 armies and soon an internal counter-revolution. The economy and working class were decimated. As a result, the democratic organs of the working class atrophied. In the resulting power vacuum, a new bureaucracy emerged to control the country and feather its own nest. Soon Russia was forced to compete with the West and as such had to accumulate capital. This new bureaucracy headed by Stalin became a new ruling class that held down workers' living standards to accumulate capital. Its rule was based on the dispossession of the working class from political power and collective ownership of the means of production. Stalin and his followers created theories to justify this counter-revolution as a continuation of the revolution. In reality however, their rule was based on the destruction of that revolution and as such was anti-socialist. They created bureaucratic state capitalism. They rejected two cardinal principles of Marxism and Leninism: internationalism and democracy. Instead of trying to spread the revolution they talked of ``Socialism in One Country''. This meant downplaying revolution and instead making government to government alliances with the capitalist West. In 1936--37, Stalin sabotaged the Spanish revolution by saying it must stay at the ``bourgeois stage'' even though workers had already begun to take power. This recreated the Menshevik (reformists in Russia) theory that in 1917 Russia was not ripe for socialist revolution which justified their support for the capitalists. This theory justified Stalin's attempt to maintain an alliance with the British and French governments which were afraid of the workers revolution in Spain.

By the late 20s, the Communist Parties had become instruments of Russian foreign policy. They therefore veered ``left'' and ``right'' according to Stalin's perception of the needs of the Russian state capitalist government. From 1929 to 1934 the CPs were ordered to be ultra-left, rejecting united fronts with other workers. This helped to lead to Hitler's rise to power in Germany. In panic Stalin then shifted the CPs back to a reformist strategy, the Popular Front. He ordered them to support ``peaceful'' capitalist parties and governments who might ally with Stalin against Hitler. The idea of revolution was thrown out the window. Later in 1939, there was switch back to a verbal leftist line of opposition to Western imperialism in the Western democracies when Stalin allied himself with Hitler.

Stalinism, though it claimed to be a continuation of Leninism, was the opposite: just another form of socialism from above. It was a bureaucratic dictatorship OVER the proletariat. Today, most of the Left that claims to be Marxist-Leninist is actually Stalinist and thus rejects the central features of Marxism, the revolutionary self-emancipation of the working class and internationalism. Because Stalin's strategy shifted over his nearly 30 years in power, the politics of these groups often differ according to which period of Stalinism they consider to be the real Stalinism. The official Communist Parties supported Russia through the Brezhnev era. They supported ``peaceful coexistence'' with Western capitalism, which fit in very well with their shift to reformism. It was another form of class collaboration that all reformism is based on. At the other extreme, some Stalinist groups look to Stalin's Third Period. They are ultra-left and generally reject united fronts with other workers. The Revolutionary Communist Party is one of these (though they often combine 3rd period ultra-leftism with the Popular Frontism). Another is the Progressive Labor Party that rejects even critical support to the nationalism of the oppressed.

\paragraph{Trotsky vs. Stalin}

In response to Stalin's perversion of Lenin and Marx's heritage, Leon Trotsky tried to preserve the real content of that tradition. He supported internationalism and democracy. Instead of ``socialism in one country'', he called for a continuation of the revolutionary policy of the early Communist International (Comintern). He opposed selling out workers' revolution in China in 1925--27 and Spain in 1936 to gain Western support as Stalin did. Instead of flipping from ultra-left rejection of a united front with other workers to reformist popular fronts with the capitalists, he urged a consistent Leninist policy of independent revolutionary parties forming united fronts with other workers' parties around immediate demands such as fighting fascism. He condemned Stalin's bureaucratic rule and called for a return to soviet democracy. Instead of subordinating the revolutionary party to bourgeois national liberation movements, he urged workers to lead the movements and take them forward to socialist revolution which he called the theory of permanent revolution. He based this in part on the practice of the Bolsheviks in Russia. For his efforts, Trotsky was expelled from Russia, hounded, isolated and finally assassinated by an agent of Stalin. However, his writings and attempt to build revolutionary organizations that continued the real Marxist and Leninist tradition were what saved later revolutionaries years of work in reviving it. Before his death, Trotsky helped form the Fourth International of very small groups around the world. Firebrand comes from that tradition.

\paragraph{Trotskyism After Trotsky}

After World War II, the Trotskyist movement split over new political developments, primarily the Russian takeover of Eastern Europe and the imposition of the Russian-style economy there. Trotsky had continued to call Russia a ``degenerated workers' state'' until his death in 1940. By this he meant that in spite of the rise of the bureaucracy, the workers' state of 1917 had never been overthrown. In spite of all his criticism of the Russian state, he felt the economy was still in a transitional stage between capitalism and socialism and was thus progressive over capitalism. By 1940, Trotsky was wrong since the workers had lost all power in the state to the new ruling class and a form of capitalism had been restored. At least there was a historical case to be made for Trotsky's position since there had been a workers' revolution in Russia in 1917. But in Eastern Europe no such revolution had ever happened. Yet the same form of economy now existed there that existed in Russia. How was this possible? If Trotsky was right that Russia was still a workers' state, the Eastern European states must be workers' states as well. If so, then workers' states could be formed without workers' revolutions! Stalinism, instead of being a counter-revolutionary force as Trotsky had argued, would become a revolutionary one since invading armies could create ``workers' states'' (the first step toward socialism). If this was true, there would be very little left of socialism as the self-emancipation of the working class. Ironically, Trotsky's description of Russia threatened to undermine his attempt to maintain the Marxist and Leninist heritage.

Many of those who stuck with Trotsky's description followed out the non-Marxist logic of it to one degree or another. Some totally collapsed into Stalinism since it was ``proved'' that Stalinism was now a revolutionary force when it carried ``revolution'' into Eastern Europe. Some of these ``Trotskyists'' rejoined the Communist Parties. A less extreme version of this viewpoint is held by the Workers' World Party which opposed the Hungarian Revolution in 1956 and Solidarity in Poland in 1980. An offshoot of Workers' World, The Party of Socialism and Liberation (PSL) was formed in 2004. This group dominates the anti-war coalition ANSWER. One way to attempt to avoid the logic of the ``workers' state'' position was to take the illogical position that Russia was still a workers' state but Eastern Europe was state capitalist. This view was held by Lutte Ouviere in France and its co-thinkers, Spark in the U.S.

For most Trotskyists, the non-Marxist logic of the ``workers' state'' analysis had a destructive effect on their politics, moving them away from focusing on workers' self-emancipation. They often began to look to other forces to bring socialism. In Europe in the 60's, some looked to students as the new group that could bring revolution. Many looked to third world nationalists who claimed to be Marxists. Today, the Freedom Socialists and other Trotskyists see Cuba as a workers' state on the road to socialism. The FSP until recently even saw China as so revolutionary that it did not need another revolution! Today, the FSP is more critical of China but still calls it a workers' state.

Most Trotskyist groups supported Russia and other state capitalist countries against the U.S. during the Cold War. They saw them as progressive. This led the Spartacist League to say ``Hail the Red Army'' about Russia's invasion of Afghanistan in 1979.

Of course, none of this is to deny that as Lenin said ``the main enemy is at home''. We must be uncompromising in our opposition to U.S. imperialism. We focus on the U.S., not because it is less progressive than other states, but because we live here and our main goal is to break workers away from support for ``their own'' government. Also, the fact that the U.S. is the major imperialist power adds urgency to the focus on opposing U.S. imperialism. However, in opposing the U.S., we need to be clear that we don't support other imperialist powers. We do of course support struggles for national self-determination against imperialism. This clarity is necessary in building a socialist movement which can defeat capitalism. We need to make sure that workers know that the movement we are building has nothing in common with the anti-worker tyrannies of China, North Korea or the former USSR. Trotskyists who give support to these tyrannies cannot give a clear definition of socialism.

\paragraph{The International Socialist Tendency}

The I.S. Tendency came out of the post-WWII Trotskyist tradition. However, it dealt with the problems created by the Russian conquest of Eastern Europe in an entirely different way than the rest of the Trotskyists. It reaffirmed the fundamental principles of Marx, Lenin and Trotsky at the expense of the descriptive, temporary aspects of their thought. The invasion of Eastern Europe and the creation of Russian style regimes there convinced our tendency that the current Russian state itself could no longer be a workers' state. We started the other direction from most Trotskyists with the undisputed facts and basic Marxist theory. There was no workers' revolution in Eastern Europe, therefore, there was no workers' state in Eastern Europe. The identical economy in Russia could not therefore be controlled by a workers' state. This led to a reexamination of Russian history and a determination that Stalin had led a counter-revolution by 1929.

Stalin had to massacre thousands of old Bolsheviks to consolidate his rule. Since then, the drive of the economy had been profit and accumulation, not human need. This theory of ``Bureaucratic State Capitalism'' was a foundation of our tendency. It allowed us to be clear about what socialism is and to oppose both major imperialist powers instead of apologizing for either. It also allowed us to keep clearly in mind that socialism can only be created by workers' revolution, not invading armies or other forces.

\paragraph{The Permanent Arms Economy}

A second major contribution of our tendency was the ``Permanent Arms Economy''. Trotsky and indeed most Marxists had expected a return to the depression after World War II. When this didn't happen, some former leftists concluded that capitalism had overcome its contradictions and revolution would never happen. Others, including many Trotskyists, tried to stick to the letter of Trotsky's predictions and claim that the post-war boom did not really exist. This denial of reality often led to predictions of immediate revolution and frantic activity followed by burnout. On the other hand, the Permanent Arms Economy explained why arms spending led to the post war boom but also why the boom would end and revolution would be on the agenda again. It allowed steady sane preparatory work in building a revolutionary tendency without on the one hand burning out or on the other giving up on revolution at all. It was based on Marx's analysis of the basic contradictions of capitalism and looking at the reality of the boom squarely in the face.

\paragraph{The Union Bureaucracy}

Another contribution was the analysis of the trade union bureaucracy. Marxists before and especially Trotskyists had denounced the conservative role of the bureaucracy. But our tendency clearly explained the inevitability of this conservative role based on its social/economic position. This clarified our union practice by showing that rank and file independence from the bureaucracy was always necessary, even when union full timers seemed to be doing the right things. It also clarified and reinforced our emphasis on rank-and-file working class self-emancipation and the need for a revolutionary party to intervene in the unions and other struggles as an independent force. This understanding of the role of the bureaucracy in the West was made easier by the theory of the Eastern bureaucracy.

\paragraph{Party And Class}

Yet another contribution was the clarification of the relationship of the revolutionary party to the working class. Under the influence of Stalinism, most in the left conceived of the party as an outside group ordering the workers around. Even Trotskyism, with its proclamation of a new international made up of small groups and no real parties in 1938, distorted the meaning of the Leninist party. Our tendency is one of the few that does not consider itself to be the revolutionary vanguard party and yet still feels that such a party is necessary. From the Communist Party to the Revolutionary Communist Party to even some Trotskyists most groups see themselves as the party that will lead the revolution. We do not see ourselves as that party. Instead, we are a tendency that wants to be part of the process of building such a party. This is not just modesty on our part. For Lenin, the vanguard party was not just a small group outside of or on the fringes of the working class. It was a large section of the working class, the most militant, class-conscious workers who actually lead struggles. None of the left groups today has come anywhere near this point. By proclaiming themselves the vanguard, they misconstrue what we are trying to build. They also show excessive idealism by overemphasizing the role of ``correct'' ideas as defining the vanguard instead of the actual activity of the vanguard in leading struggle.

Such distortions also downplay the role of the working class in its own liberation by saying it can be led by a small group of outsiders with the ``right'' ideas. Another important part of this development of the theory of party and class was the clarification that the party was not the same as the state. For both Stalinists and Social Democrats, the socialist party and the state are virtually identical. For Lenin the Soviets, or workers' councils, which included the whole working class, were the basis of the workers' state. The party would have influence in the state by winning support in those councils. The workers could and did elect different parties to the soviet.

\paragraph{Deflected Permanent Revolution}

Still another contribution was the theory of ``deflected permanent revolution''. Trotsky's analysis of the Russian Revolutions of 1905 and 1917 rejected the reformist and later Stalinist theory of stages. He did not believe there had to be a full stage of capitalist class rule before the working class could come to power. In Russia the working class was strong and concentrated while the bourgeoisie was weak and afraid that even a bourgeois revolution would get out of control and lead to attacks on its property. Because of this, the proletariat could pick up the leadership of the bourgeois revolution and carry it forward to a seizure of power by the workers. Of course, for the workers to stay in power in a mainly peasant country, the revolution would have to spread to other countries. Thus, the revolution must become permanent, or continuous, in two senses: passing from the bourgeois stage over to the socialist stage and from one country to another. Since the bourgeoisie was too weak and timid to lead the revolution, Trotsky saw no other force that could lead it besides the workers. Trotsky's theory was proven in the Russian Revolution. Those socialists who accepted the stage theory actually supported the bourgeoisie in counter-revolution against the Bolshevik-led Soviets from 1918 on. As we've seen, Stalin's stage theory led to defeats in China and Spain.

Yet if Trotsky's theory was correct, as a guide to action, his prediction that there was no other force that could lead a revolution against imperialism or feudalism besides the proletariat was shown to be wrong. After World War II virtually the whole colonial world won independence without a proletarian revolution. The leadership in many cases came from urban professionals who could benefit from state independence especially if the state seized much of the economy and they managed the state. Often, this revolutionary leadership became a new state capitalist ruling class (especially in China, Cuba, Viet Nam, and North Korea) but to a lesser extent even in the mixed state/private capitalist economies of other countries.

For Trotskyists, this presented a problem. Since Trotsky implied that only a proletarian revolution could win independence from imperialism, either these revolutions were proletarian or Trotsky was wrong. (Or these countries weren't really independent since their economies were still controlled by the old colonial power. But this variation would not fit the state capitalist countries since the old imperialists were expelled.) Instead of modifying the theory to fit the facts, most Trotskyists modified the facts to fit the theory. They said the state capitalist rulers had led a workers' revolution, when it was clear the workers played little role in many of these revolutions, at least little independent role. Their tendency to see state capitalism as really a form of ``workers' state'' was reinforced by their similar analysis of Eastern Europe.

This new theory reinforced the illusion that socialism could be a product of something else besides working-class action, in this case nationalist guerilla movements. Today, some Trotskyists carry this so far as to identify Castro with Lenin, seeing Cuba as almost a full workers' state rather than a ``deformed workers state''. The Freedom Socialist Party (FSP) is almost uncritical of Cuba even on the issue of gay oppression, which is otherwise one of its main issues.

Our tendency dealt with post-war independence in a different way. We accepted the core of Trotsky's strategy and rejected the stages theory. The task of socialists was as always to build working class independence and revolutionary parties that could lead successful workers' revolutions and begin the process of international permanent revolution. Instead of distorting facts to fit theory, we recognized the fact that independence had been won without working class leadership and therefore the theory needed modification. The theory of ``deflected permanent revolution'' did this by examining the factors that led to the failure of the working class to take the lead in the anti-colonial revolution. The primary factor was a political failure, the lack of strong revolutionary socialist parties and instead the dominance of Stalinism and the model of Russian state capitalism as a road to economic independence. Another important factor was that imperialism no longer needed formal ownership of colonies as much as it had before the war. When the working class was not politically prepared to fight for power, and the bourgeoisie was too timid, and imperialism's grip was weakening, a new incipient state capitalist class which could benefit from independence could lead the movement for independence.

There was still no historical necessity for this state capitalist stage. It was a question of whether the workers were politically prepared or not. Therefore, the conclusion of the theory was the same as Trotsky's. If anything, the need to jettison the stages theory and organize for workers power was even stronger. The post war changes reinforced this task. The working class was larger in all these countries and less of the old pre-capitalist structures existed. More and more the revolution needed in underdeveloped countries was the same as in the industrialized ones: a working class socialist one. This attitude differentiated us from even most Trotskyists who looked at nationalist guerrilla leaders as the bringers of ``workers' states''.

Of course, none of this is to deny the overwhelming importance of opposing imperialism, especially U.S. imperialism since we live in the U.S. We may not think that Cuba or Vietnam is socialist but their right to be independent from imperialism is absolute. We must first of all emphasize this and only then raise our criticisms of those regimes.

This is a brief summary of the historical basis of the I.S. Tendency. Only the tradition of Lenin and Trotsky preserved the essential vision of Marx of working-class self-emancipation and internationalism as the core of socialism. After the death of Trotsky most of his followers veered back toward Stalinism to one degree or another, thereby compromising his vision of socialism as based on workers' democracy. They also distorted the Leninist idea of the revolutionary vanguard party by overemphasizing a ``correct program'' to the exclusion of the working-class basis of the party. Both of these problems undermined their commitment to a working-class orientation as opposed to the labor bureaucracy or other groups.

\paragraph{Weaknesses of the International Socialist Tendency}

In spite of the important contributions that the IST made to the revolutionary movement, it had several weaknesses. The Tendency was organized from London in a top-down manner. The International Socialist Organization, the predecessor of Red Spark, was expelled from the Tendency in 2001 over differences in perspective. From that point on the I.S.O. applied the general politics of the Tendency adapting those politics to the U.S. situation. The top-down nature of the IST by 2001 reflected a degeneration of its view of the revolutionary party. The SWP-GB, the main organization in the IST, was internally less democratic than the Leninist tradition required. Unfortunately, the influence of the IST led to similar weaknesses in the ISO. Besides this, the SWP tended to be too dismissive and hostile to ideologies of liberation such as Feminism. It also veered between being sectarian toward movements and being overly accommodating.

\paragraph{Red Spark and the ISO}

Red Spark was formed out of the dissolution of the ISO in 2025. There were several overlapping issues that caused the dissolution. One was the influence of the aforementioned top-down structure of the IST on the ISO. Though the ISO broke from the IST it never fully broke from the IST's methods. Some comrades went so far as to describe relationships between the leaders and the rank and file as abusive. Another issue was that many members of the ISO were moving in an opportunistic direction in response to the success of DSA from 2015 onward and the influence of the Bernie Sanders campaign. In order to relate to the ``New Socialist Movement'', a wing of the ISO called ``Socialist Tide'' called for support of the Sanders campaign. In doing so, they jettisoned a fundamental principle of Marxism, the need for the working class to be politically independent from the capitalist class and its parties. There was also frequent mishandling of the oppression within the group.

The final straw that caused the dissolution was the revelation of a mishandled rape accusation in 2013. It turned out that a new member of the ISO leadership that took over in February 2019 had been accused of rape. After details about how the case was handled, previously known only by some members of the 2012 steering committee and a couple other bodies of the organization were made public, the integrity of the process and the ruling were severely called into question by many members. For not a few comrades, this could be accurately described as a coverup. Combined with the other political problems in the ISO, this led a majority of members and ex-members to vote for dissolution in March 2019.

\paragraph{Red Spark Going Forward}

We believe that clarity on the issues that have divided the Left are important to understand. In each historical period, it was crucial for revolutionary continuity to be on the correct side of these issues. However, the immediate practical importance of particular theories changes over time. Today, key dividing lines on the revolutionary left are different than during the Cold War. An exact analysis of Stalinism is still important in the long run but is not enough to form different groups now. This was recognized by the ISO when it removed adherence to the State Capitalist analysis as a fundamental point in its Where We Stand a few years ago.

Today the revolutionary and would-be revolutionary Left is primarily divided over 5 key questions: 1) Self-emancipation of the working-class, or elitist rule 2) Should Marxists support Democratic Party candidates? 3) Should Marxists support workers in their struggles against ruling classes in every country of the world---or should they instead oppose workers' movements in countries that are enemies of the U.S.? 4) Should Marxists combine the fight against exploitation with the fight against oppression? Or should Marxists take a class reductionist position? 5) Should Marxists support independent rank and file organizing in the unions?

\begin{enumerate}
\item Self-emancipation of the working class. As this article points out, the Left since the late 20's has been divided between those who support Socialism from Above, whether in the form of Stalinism or Social Democracy and those who actually support grass roots working class democratic control of society. Those who see Stalin's Russia or Mao's China for example as socialist societies are actually re-writing socialism into its opposite whether they intend to or not. They are supporting bureaucratic rule rather than workers' power. We support workers' power and are on the Trotskyist side of the Trotsky-Stalin split.
\item As noted above, Red Spark has been clear that it sees support for Democratic candidates, even Sanders, AOC etc. as an abandonment of the Marxist principle of working-class independence. Those who claim to be revolutionaries while involving themselves in Democratic campaigns are confusing people on what the revolutionary Marxist position really is.
\item Some on the Left defend oppressive regimes abroad because they are supposedly enemies of U.S. imperialism---Putin, Assad, Iran etc. These groups (PSL, WW etc.) abandon the principle of self-emancipation of the working class. On the opposite side, some on the Left support candidates who do not staunchly oppose U.S. imperialism. As German revolutionary Karl Liebknecht said ``The main enemy is at home''. Our first responsibility is to oppose U.S. imperialism. However, following Liebknecht, we understand that while the MAIN enemy is at home, the U.S. ruling class is not the ONLY enemy. The ruling class enemies of the workers of Iran, Syria, Russia, Venezuela etc. are also our enemies. We are internationalists. As such, we oppose the enemies of workers everywhere.
\item Some in DSA and other groups adopt class reductionism and put so much emphasis on supporting universal demands such as Medicare for All that they ignore Lenin's approach of supporting special demands of the oppressed. On the other hand, some on the Left put so much emphasis on special oppression that they nearly ignore the need to also support universal demands and the fight against exploitation. In the long-standing Marxist tradition, we seek to combine the fight against oppression and exploitation. We fight against all forms of oppression because oppression is one of the most disgusting and revolting aspects of capitalism. We also oppose oppression because a revolutionary movement that does not take on the fight against oppression will never create enough working-class unity to succeed. The fight against exploitation cannot succeed without a simultaneous fight against oppression. Likewise, oppression cannot be destroyed within the confines of the capitalist system. Only a successful victory against exploitation, i.e. the elimination of the capitalist system can lay the basis for smashing all forms of oppression. We oppose both class reductionism and separatism.
\item Red Spark supports unions but understands that the social position of the union bureaucracy often leads it to hold back struggle. We support independent rank and file organization within the unions to counter this tendency. We fight for democracy in the unions as well as in all social movements.
\end{enumerate}

We see these five points as the fundamental stands that revolutionary socialists must unite around in order to be effective today. We work with anyone, socialist or not, in coalitions around particular issues. We seek even closer unity with other revolutionaries who agree on these principles today. Red Spark today includes former ISO members and Trotskyists who have supported the idea that Russia was a degenerated workers' state, not State Capitalist. We seek to expand our collaboration with others on the revolutionary Left, especially those that support these principles.

Red Spark sees itself as part of the process of cohering larger organizations of revolutionaries. The long-term goal is to be part of the process of building an actual working class revolutionary vanguard party like the Bolshevik Party that led the Russian Revolution. The steps along that way are so far uncharted. Please join us in this effort!
